\subsection{The exact global freedom in a class-two lift}
\label{sec:ghost-h1-weight2}

\begin{result}{C.1.18}{Proved}{Continuous first cohomology in weight two}
\label{prop:ghost-h1-weight2}
There are isomorphisms
\[
 H^1_{\rm cts}(G_{\mathbb Q},\mathbb Z_\ell(2))
 \simeq
 \begin{cases}
  \mathbb Z/8,&\ell=2,\\
  \mathbb Z/3,&\ell=3,\\
  0,&\ell\ge5,
 \end{cases}
 \qquad
 H^1_{\rm cts}(G_{\mathbb Q},\widehat{\mathbb Z}(2))
       \simeq\mathbb Z/24.
\]
The latter group is generated by the integral continuous
cocycle \(\eta(g)=(\chi(g)^2-1)/24\).
\end{result}
\begin{proof}
We retain the coefficient transition maps throughout.
The universal coefficient sequence
\cite[IV.2.5]{GhostWeibelKbook} is
\[
 0\longrightarrow K_3(\mathbb Q)/\ell^n
 \longrightarrow K_3(\mathbb Q;\mathbb Z/\ell^n)
 \longrightarrow K_2(\mathbb Q)[\ell^n]
 \longrightarrow0.
\]
Its transition from exponent \(n+1\) to \(n\) is
reduction on the first term and multiplication by \(\ell\)
on the last term. Indeed, form the coefficient spectrum
as the cofiber of multiplication by \(\ell^n\) on the
\(K\)-theory spectrum. The morphism of cofiber sequences
is multiplication by \(\ell\) on the source spectrum
and the identity on the target spectrum. The homotopy
long exact sequence gives precisely these two maps.

The tame-symbol exact sequence
\cite[III.5.2.2, III.6.5.1]{GhostWeibelKbook} is
\[
 0\longrightarrow\mathbb Z/2\longrightarrow K_2(\mathbb Q)
 \longrightarrow\bigoplus_q\mathbb F_q^\times
 \longrightarrow0.
\]
Consequently \(T_\ell K_2(\mathbb Q)=0\).
In fact, a compatible sequence \(x_n\in K_2(\mathbb Q)[\ell^n]\)
with \(x_n=\ell x_{n+1}\) has zero image in each finite
group \(\mathbb F_q^\times\): iterating the compatibility
more times than the exponent of its \(\ell\)-primary
part kills that coordinate. Each \(x_n\) therefore lies
in the displayed kernel \(\mathbb Z/2\), in which the
same iteration kills the entire sequence. Since the
finite groups \(K_3(\mathbb Q)/\ell^n\) have zero
derived inverse limit, the coefficient sequence now gives
\[
 \varprojlim_n K_3(\mathbb Q;\mathbb Z/\ell^n)
       =K_3(\mathbb Q)\otimes\mathbb Z_\ell.
\]

The natural finite-coefficient motivic edge sequence
\cite[VI.4.3.1]{GhostWeibelKbook} has the exact tail
\[
 K_3^M(\mathbb Q)/\ell^n
 \longrightarrow K_3(\mathbb Q;\mathbb Z/\ell^n)
 \longrightarrow H^1(G_{\mathbb Q},\mathbb Z/\ell^n(2))
 \longrightarrow0.
\]
The integral edge map is injective, and its groups are
\[
 K_3^M(\mathbb Q)=\mathbb Z/2,
 \qquad K_3(\mathbb Q)=\mathbb Z/48;
\]
see \cite[VI.4.3.2, VI.5.2.1, VI.5.3]{GhostWeibelKbook}.
Thus its image is the unique order-two subgroup.
Every finite-coefficient kernel in this tail is finite,
so its inverse system has zero derived inverse limit.
For \(\ell=2\) and \(n\ge4\), the Milnor order-two
element remains nonzero in
\(K_3(\mathbb Q)/2^n\), hence in the coefficient group;
these kernels are compatibly that order-two subgroup.
For odd \(\ell\) they are zero. It follows that
\[
 \varprojlim_n H^1(G_{\mathbb Q},\mathbb Z/\ell^n(2))
 =\bigl(K_3(\mathbb Q)/K_3^M(\mathbb Q)\bigr)
                       \otimes\mathbb Z_\ell
 = (\mathbb Z/24)\otimes\mathbb Z_\ell.
\]
This is continuous integral first cohomology: the
cochain complexes reduce surjectively in every degree,
and the only comparison error in degree one is
\(\varprojlim{}^1 H^0(G_{\mathbb Q},\mathbb Z/\ell^n(2))\),
which vanishes because all its terms are finite.

For the explicit generator put \(T=\mathbb Z_\ell(2)\)
and \(V=\mathbb Q_\ell(2)\). A continuous cocycle into
\(V\) has bounded image, since its domain is compact.
Multiplying it by a power of \(\ell\) gives an integral
cocycle. The finite integral first cohomology just
computed therefore proves \(H^1_{\rm cts}(G_{\mathbb Q},V)=0\).
Also \(H^0(G_{\mathbb Q},V)=0\). The connecting map
for \(0\to T\to V\to V/T\to0\) is consequently an
isomorphism from \((V/T)^{G_{\mathbb Q}}\) to \(H^1(G_{\mathbb Q},T)\).
Surjectivity of the cyclotomic character identifies
these invariants with the elements killed by every
\(u^2-1\), for \(u\in\mathbb Z_\ell^\times\).
The ideal generated by these differences is
\(8\mathbb Z_2\), \(3\mathbb Z_3\), or \(\mathbb Z_\ell\),
respectively: at two use \(u=3\), at three use \(u=2\),
and for \(\ell\ge5\) choose a nonzero residue other
than \(\pm1\). Thus the class of \(1/24\) generates
the invariant subgroup in every component, and its
connecting cocycle is exactly \(\eta\).
Continuous cochains with product coefficients are the
products of the component cochains; their cohomology
is the product as well. This proves the profinite
assertion and the stated generator.
\end{proof}

\begin{result}{C.1.19}{Proved}{Exactly twenty-four class-two section classes}
\label{cor:ghost-class2-fibre24}
For every ordinary completed pair \((A,B)\), the fibre
of geometric-conjugacy classes of continuous sections
of the actual class-two arithmetic quotient above
\((A,B)\) is a simply transitive \(\mathbb Z/24\)-set.
For the geometrically pro-\(\ell\) quotient, the
corresponding fibre has eight elements at two, three
at three, and one at every \(\ell\ge5\).

More explicitly, choose literal abelian cocycles \(a,b\)
and one lift with central coordinate \(h\). Representatives
for all twenty-four classes are
\[
       y^{b(g)}x^{a(g)}C^{h(g)+k\eta(g)}s_0(g),
                 \qquad k\in\mathbb Z/24.
\]
\end{result}
\begin{proof}
Existence follows from Corollary~\ref{cor:ghost-class2-lift}.
With \(a,b\) fixed literally, the difference of two
central coordinates is a continuous cocycle in
\(\widehat{\mathbb Z}(2)\). Central geometric conjugation
changes it by a continuous coboundary, so its orbits
form a torsor under the first cohomology just computed.

This accounts for full geometric conjugation, not merely
for central conjugation. First, any representative of
the same abelian cohomology class can be conjugated by
a lift of an abelian geometric element so that its two
abelian cocycles become exactly \(a,b\). Next, a
geometric element conjugating two such normalized
sections has an abelian coordinate fixed by
\(G_{\mathbb Q}\). That invariant group is zero:
\[
 H^0(G_{\mathbb Q},\widehat{\mathbb Z}(1)^2)=0,
\]
because each component is torsion free and some
cyclotomic scalar differs from one. The conjugator
therefore lies in the central class-two layer.
No further identifications are made. The same argument
componentwise proves all the pro-\(\ell\) assertions.
\end{proof}
